\scanpage{75}
% archive foliation: 77
identifiant les pts des centres des disques
qui se correspondent, et en ponctuant
par les centres des autres disques.
De plus, on obtient, \deleted{avec} pour chaque
point singulier $x$ de $X'''$, correspondant
à une paire $\{x_1,x_2\}$ de points de
$X''$ (centres des deux disques associés),
grâce à l'antiisomorphisme de ces
disques pointés, un élément \deleted{de}\inserted{non nul (i.e. une base)}
de
\[
T_{X'',x_1}\otimes_{\mathbb C}T_{X'',x_2}
\]
i.e. un élément de
\[
T^*_{X'',x_1}\wedge_{\mathbb C^*}T^*_{X'',x_2}.
\]

On trouve ainsi un foncteur de

a) catégorie \inserted{(groupoïde)} des courbes hol. cptes à bord \inserted{$X$} de
type $g,\nu$, avec découpage holomorphe
(avec comme flèches les iso)

vers

b) Catégorie des courbes \inserted{$\mathbb C$-alg} stables de type $(g,\nu)$
$X'''$, munies \inserted{1°)} d'une ``structure de propagation
\unsure{unnodale|nodale}'' en ses pts singuliers, et 2°) d'un
\deleted{structure de \unsure{parallélismes}}\inserted{syst.} admissible de
disques autour des pts marqués.
